\chapter{Considerações Finais}

O Projeto Pedagógico do Curso de \ppccurso da Universidade Federal de Uberlândia reflete o compromisso da instituição com a formação de profissionais éticos, tecnicamente qualificados, inovadores e socialmente responsáveis. Estruturado em consonância com a Resolução CNE/CES nº 5/2016 (DCN Computação), com as diretrizes do \textit{Computing Curricula 2020} (CC2020/ACM) e com o Projeto Institucional de Desenvolvimento e Expansão (PIDE) da UFU, este PPC apresenta uma proposta formativa fundamentada em competências, interdisciplinaridade, integração teoria-prática e valorização da diversidade humana e cultural.

A matriz curricular, os objetivos do curso, o perfil do egresso, as estratégias de ensino e os mecanismos de avaliação foram cuidadosamente definidos para garantir uma formação alinhada aos desafios tecnológicos contemporâneos e às demandas da sociedade. O curso busca oferecer ao(à) estudante um percurso formativo flexível, crítico, inovador e capaz de dialogar com a complexidade do mundo do trabalho, da ciência e da vida em sociedade.

Este documento expressa também a valorização do diálogo entre discentes, docentes, técnicos(as) e gestores(as), bem como o papel estratégico da universidade pública no desenvolvimento regional e nacional. O curso de \ppccurso da FEELT/UFU reafirma, assim, sua missão de formar engenheiros(as) comprometidos com a excelência acadêmica, a justiça social e a construção de um futuro sustentável.
